\scanpage{62}
tous nos beaux énoncés devraient être valables, non seulement pour un corps de base $K$ ext. de type fini de $\mathbb{Q}$, mais dès que $K$ est ext. de type fini d'un corps cyclotomique (pas néc. fini sur $\mathbb{Q}$).

Nous pourrions définir les \uline{courbes de Poincaré} sur un corps alg. clos $\bar K$ de car. $0$ comme étant les courbes isomorphes à des rev. universels de courbes alg. anabéliennes sur $\bar K$ (deux courbes anabéliennes $\unsure{=}\ \bar U,\bar V$ sur $\bar K$ définissent des rev. de Poincaré isomorphes, ssi $\exists$ un revêtement fini étale de l'une, isomorphe à un rev. fini étale de l'autre\ldots).

Étant donné une courbe de Poincaré $\widetilde{\unsure{\mathcal U|U}}$ sur $\bar K$, \deleted{telle} on définit canoniquement sa complétée\inserted{ $\widehat{\unsure{\mathcal U|U}}$}, de telle façon que $\forall$ autom de $(\widetilde{\unsure{\mathcal U|U}},\bar K)$ (pas néc. induisant l'identité sur $\bar K$) induise un autom. de $\widehat{\unsure{\mathcal U|U}}$.
Ceci posé :
\margin{61}
